\documentclass[../../main.tex]{subfiles}
\begin{document}

\begin{center}
  {\fontsize{48pt}{58pt}\selectfont \textbf{2010}\normalsize\textbf{年度}}
\end{center}

\begin{table}[htb]
  \centering
  \caption{出題テーマと難易度}
  \begin{tabular}{|c|p{6.2cm}|c|c|c|}
    \hline
    \textbf{問題番号} & \textbf{テーマ} & \textbf{思考量} & \textbf{計算量} & \textbf{総合} \\
    \hline
    第1問 & 直方体の回転でできる立体の体積 & 易 & 普 & 易 \\
    \hline
    第2問 & 定積分の評価と対数の不等式 & 普 & 普 & 普 \\
    \hline
    第3問 & 二つの箱の玉の移動と確率漸化式 & 普 & 普 & 普 \\
    \hline
    第4問 & 曲線上の二点と面積 & 難 & 易 & 普 \\
    \hline
    第5問 & 円周上を動く三点と直角二等辺三角形 & 普 & 易 & 普 \\
    \hline
    第6問 & 四面体の切り口の面積 & 普 & 難 & 難 \\
    \hline
  \end{tabular}
\end{table}

\begin{table}[htb]
  \centering
  \caption{解答}
  \renewcommand{\arraystretch}{1.4}
  \begin{tabular}{|c|c|p{9cm}|}
    \hline
    \textbf{問題番号} & \textbf{小問} & \textbf{答え} \\
    \hline
    第1問 & (1) & $V=b\left\{ac+\dfrac{\pi}{4}(a^2+c^2)\right\}$ \\
    \cline{2-3}
    & (2) & $0<V<\dfrac{\pi}{27}$ \\
    \hline
    第2問 & (1)・(2) & 定積分と対数の上下界を示す（証明問題） \\
    \hline
    第3問 & (1) & $P_m(x)$を$P_{m-1}$で表す漸化式（解答参照） \\
    \cline{2-3}
    & (2) & $P_{2n}(10)=\dfrac{1}{3}\left\{1-\left(\dfrac{1}{4}\right)^n\right\}$ \\
    \cline{2-3}
    & (3) & $P_{4n}(6)=\dfrac{1}{5}\left\{1-\left(\dfrac{1}{16}\right)^n\right\}$ \\
    \hline
    第4問 & (1) & $\triangle OP_1H_1$と$\triangle OP_2H_2$は等積（証明問題） \\
    \cline{2-3}
    & (2) & $2\log\dfrac{y_2}{y_1}$ \\
    \hline
    第5問 & -- & $m=4$のとき$t=\dfrac{\pi}{6},\dfrac{\pi}{2},
    \dfrac{5\pi}{6},\dfrac{7\pi}{6},\dfrac{3\pi}{2},\dfrac{11\pi}{6}$．
    $m=8$のとき$t=\dfrac{\pi}{2},\dfrac{3\pi}{2}$ \\
    \hline
    第6問 & (1) & $\overrightarrow{OH}=\dfrac{5}{9}\overrightarrow{OA}
    -\dfrac{1}{3}\overrightarrow{OB}$ \\
    \cline{2-3}
    & (2) & $0<t\le\dfrac{2}{9}$なら$S(t)=3\sqrt{6}t^2$，
    $\dfrac{2}{9}<t<1$なら
    $S(t)=\dfrac{12\sqrt{6}}{49}(-8t^2+9t-1)$ \\
    \cline{2-3}
    & (3) & 最大値$\dfrac{3\sqrt{6}}{8}$ \\
    \hline
  \end{tabular}
\end{table}

\paragraph{試験形態}
6問．試験時間は要確認．

\end{document}
